% =============================================================================
% Chapter 4: Data Model
% =============================================================================

\chapter{Data Model}
\label{ch:datamodel}

This chapter defines the data structures that the rest of the SRS
references. It is placed before the requirement chapters because the
requirements are stated in terms of these structures.

\section{The VarFrame}
\label{sec:varframe}

A VarFrame is the primary multidimensional data structure in \itc{}. It
stores the following attributes:

\begin{table}[H]
\centering
\caption{VarFrame attributes.}
\label{tab:varframe-attrs}
\small
\begin{tabularx}{\textwidth}{@{}llX@{}}
\toprule
\textbf{Attribute} & \textbf{Type} & \textbf{Description}\\
\midrule
\code{dims}        & \code{dict[str, Dimension]} &
Ordered mapping of dimension name to a \code{Dimension} object holding the
coordinate array, optional unit metadata, and optional non-numeric flag.\\
\code{vars}        & \code{dict[str, Variable]} &
Mapping of variable name to a \code{Variable} object wrapping the
N-dimensional NumPy array, with optional unit and description metadata.
Every variable array has the same shape, dictated by the dimension order.\\
\code{uncertainty} & \code{UncFrame \textbar\ None} &
Mirror structure for standard uncertainties; \code{None} when no
uncertainties are assigned.\\
\code{tags}        & \code{HistoryFrame \textbar\ None} &
Mirror structure tagging each value as original (\code{0}), interpolated /
computed (\code{+1}), or extrapolated (\code{-1}). Created lazily upon the
first operation that produces non-original values.\\
\code{provenance}  & \code{Provenance} &
Static origin record (source files, hashes, user, creation timestamp,
\itc{} version, operating mode).\\
\code{history}     & \code{History} &
Cumulative, ordered record of every operation applied since creation.\\
\code{axes}        & \code{AxisRegistry} &
Coordinate-frame registry containing the canonical AIAA body axis (default)
and any user-defined custom axes (Section~\ref{sec:axes}).\\
\code{coords}      & \code{CoordSystem} &
Spatial coordinate system tag: Cartesian or Polar (affects integration).\\
\code{correlation} & \code{CorrelationMatrix \textbar\ None} &
Optional correlation matrix between variables for uncertainty propagation.
\code{None} means all variables are treated as independent.\\
\bottomrule
\end{tabularx}
\end{table}

\subsection{Dimension Ordering and Shapes}

The shape of every variable array is determined by the order of \code{dims}.
If \code{dims = \{"alpha": ..., "mach": ..., "blade": ...\}} then every
variable array has shape $(N_\alpha, N_M, N_\beta)$. Ordering is preserved
across all operations.

\subsection{Missing Values}

Missing dimension combinations are represented as \code{np.nan} in the
variable arrays. \itc{} never uses sentinel values or masked arrays. The
HistoryFrame distinguishes a value that was never measured (NaN in the
variable, no entry in \code{tags}) from a value that was interpolated to
fill a gap (finite in the variable, \code{+1} in \code{tags}).

\subsection{The Dimension Object}
\label{sec:dimension-object}

Each entry in \code{dims} is a frozen dataclass:

\begin{lstlisting}
@dataclass(frozen=True)
class Dimension:
    name: str                    # e.g. "alpha"
    coords: np.ndarray           # 1-D coordinate array
    unit: str | None = None      # "deg", "m", "rpm", None
    description: str | None = None
    is_numeric: bool = True      # False for string-valued dims (e.g. blade_type)
\end{lstlisting}

The \code{unit} field is metadata only and is not used for unit conversion
unless the user calls \code{utils.units.convert()} explicitly. It is read by
the plotting system for axis labels and by export formats for header
annotations.

Non-numeric dimensions (\code{is\_numeric=False}) are permitted, with
coordinates such as \code{["A", "B"]} for blade-type sweeps. Numerical
operations on such dimensions (\code{interpolate}, \code{diff},
\code{fitmodel}, \code{integrate}) raise
\code{NonNumericDimensionError} immediately, before any computation.

\subsection{The Variable Object}
\label{sec:variable-object}

Each entry in \code{vars} is a frozen dataclass wrapping the underlying
array:

\begin{lstlisting}
@dataclass(frozen=True)
class Variable:
    name: str                       # e.g. "CT"
    values: np.ndarray              # N-D array, shape from VarFrame.dims
    unit: str | None = None         # "N", "m/s", None
    description: str | None = None
    long_name: str | None = None    # used for plot labels: "Thrust coefficient"
\end{lstlisting}

Mirroring the design of \code{Dimension}, the \code{unit} field on a
\code{Variable} is metadata only and feeds the unit-conversion utility
(\code{utils.units.convert}) on demand.

\section{The UncFrame}

An UncFrame is structurally a VarFrame restricted to standard uncertainties:
same dimensions, same variable names as its parent, with array values
representing $u(\text{var})$ following GUM JCGM~100:2008 clause~2.3.1.

It is created automatically the first time \code{db.set\_uncertainty(...)}
is called, and it is updated automatically by every \code{db.compute(...)}
that involves variables with assigned uncertainty.

\subsection{Two Uncertainty Components}
\label{sec:two-components}

For each variable, the UncFrame stores up to two arrays: a systematic
component $u_{\mathrm{sys}}$ and a random component $u_{\mathrm{rand}}$.
The distinction follows AIAA S-071A-1999~\citep{aiaaS071A}: systematic
uncertainty originates in calibration and installation biases and is fully
correlated across the points of a variable; random uncertainty originates
in measurement scatter and is independent between points.
\code{db.set\_uncertainty} assigns the systematic component by default and
accepts \code{component="random"}. The distinction matters in reductions:
averaging $N$ independent repeats reduces $u_{\mathrm{rand}}$ by
$1/\sqrt{N}$ but leaves $u_{\mathrm{sys}}$ unchanged. A single-component
model cannot represent this and silently understates the combined
uncertainty of averaged results. Both components propagate separately
through every operation and recombine as
$u = (u_{\mathrm{sys}}^2 + u_{\mathrm{rand}}^2)^{1/2}$ at reporting time
(REQ-99).

\itc{} v0.1.0 supports correlated input variables natively. The combined
standard uncertainty of a derived quantity $f$ is
\begin{equation}
u^2(f) = \sum_i \left(\frac{\partial f}{\partial x_i}\right)^2 u^2(x_i)
       + 2\sum_{i<j} \frac{\partial f}{\partial x_i}
                     \frac{\partial f}{\partial x_j}
                     \,r(x_i, x_j)\, u(x_i)\, u(x_j),
\label{eq:lpu}
\end{equation}
where $r(x_i, x_j)$ is the correlation coefficient declared via
\code{db.set\_correlation(...)}. The declared coefficient applies
identically to both uncertainty components (validated 2026-07-21,
OQ-23 in the companion log). By default all correlations are zero and
Equation~\eqref{eq:lpu} reduces to the independence form. Partial derivatives
are obtained symbolically via chain rule on the RPN expression tree
(Section~\ref{sec:rpn-engine}).

\begin{noteinfo}[Monte Carlo for discrete-branch models]
The symbolic propagation in Equation~\eqref{eq:lpu} assumes the derived
quantity is differentiable with respect to its inputs. Models containing
discrete branches (piecewise definitions, conditional logic that selects different equations based on input values) break this assumption. For
such cases \itc{} provides Monte Carlo propagation
(\code{db.compute(..., method="mcm")}) following GUM
S1~\citep{gum2008s1}. Continuous nonlinearities are handled correctly by
the default symbolic propagation and do not require Monte Carlo.
\end{noteinfo}

\section{The HistoryFrame}

The HistoryFrame tracks the origin of every value in every variable. For
each variable in the parent VarFrame, it stores an integer array of
identical shape with values:

\begin{itemize}[noitemsep,topsep=2pt]
  \item \code{0}: value was loaded directly from a source file (measured
    ground truth or simulated raw output).
  \item \code{+1}: value was interpolated, computed, or otherwise derived
    within the convex hull of the source data.
  \item \code{-1}: value was extrapolated outside the convex hull of the
    source data.
\end{itemize}

Operations that produce derived values are responsible for updating the
HistoryFrame accordingly. \code{compute} sets \code{+1} for newly created
variables. \code{interpolate} sets \code{+1} within the original axis range
and \code{-1} outside. \code{fill} sets \code{+1}. The reshaping rules for
\code{fitmodel}, \code{fitvalue}, and the windowed operations are stated in
the OQ-25 paragraph below. The plotter consumes
the HistoryFrame to allow visual distinction (filled markers for original,
hollow markers for interpolated, hollow with cross for extrapolated).

\paragraph{Reductions.}
Operations that collapse one or more dimensions (\code{average},
\code{integrate}, \code{statistics}, \code{combine}) follow a worst-case
rule: the resulting tag is \code{-1} if any contributing point was
\code{-1}, else \code{+1} if any was \code{+1}, else \code{0}. This
preserves the conservative meaning of the tag (a reduced result is at most
as ``original'' as its weakest contributor). For weight-based reductions
(\code{average}, \code{integrate}), the worst case is taken over the cells
that carry nonzero weight.

\paragraph{Fits, evaluation, and windows (OQ-25).}
Operations that reshape the tag grid extend the same worst-case
conservatism: \code{fitmodel} spreads the worst-case tag of a fitted line
across all of its coefficients (a coefficient is at most as ``original''
as the weakest sample that produced it); \code{fitvalue} tags an evaluated
point \code{+1} inside the recorded fit range and \code{-1} beyond it (an
extrapolation is derived, REQ-32); and windowed operations (\code{smooth},
\code{diff}) take the worst case over each moving window.

\section{Provenance and History}
\label{sec:provenance-history}

\itc{} maintains a strict separation between two concepts that are often
conflated.

\subsection{Provenance: where the data came from}

Provenance is the static, immutable origin record of a VarFrame, established
at creation time and never modified afterward.

\begin{table}[H]
\centering
\caption{Fields of the \code{Provenance} object.}
\label{tab:provenance-fields}
\small
\begin{tabularx}{\textwidth}{@{}llX@{}}
\toprule
\textbf{Field} & \textbf{Type} & \textbf{Description}\\
\midrule
\code{itaca\_version} & \code{str}     & Version of \itc{} that created this VarFrame.\\
\code{user}           & \code{str}     & \code{user@hostname}, overridable.\\
\code{created\_at}    & \code{datetime}& ISO 8601 creation timestamp.\\
\code{source\_files}  & \code{list[Path]}& Paths to input files.\\
\code{source\_hash}   & \code{str}     & SHA-256 digest of the concatenated source file contents.\\
\code{mode}           & \code{str}     & \code{"production"} or \code{"draft"}.\\
\code{version\_tag}   & \code{str}     & User-defined string (e.g.\ \code{"v1.0-raw"}).\\
\code{source\_coords} & \code{tuple \textbar\ None} & Per-file dimension coordinates captured at load time (\code{"*"} marks a dimension swept within the file); the backing record of \code{db.manifest} (OQ-20).\\
\bottomrule
\end{tabularx}
\end{table}

\subsection{History: what happened to the data}

History is the ordered, append-only sequence of operations applied to a
VarFrame since its creation. Each entry has a unique sequential index
within the VarFrame, starting from 1.

\begin{table}[H]
\centering
\caption{Fields of a single History entry.}
\label{tab:history-fields}
\small
\begin{tabularx}{\textwidth}{@{}llX@{}}
\toprule
\textbf{Field} & \textbf{Type} & \textbf{Description}\\
\midrule
\code{index}       & \code{int}      & Sequential index, starting at 1.\\
\code{operation}   & \code{str}      & Operation name and arguments.\\
\code{comment}     & \code{str}      & Optional user comment passed via the \code{comment=} keyword to any operation.\\
\code{timestamp}   & \code{datetime} & When the operation was applied.\\
\code{state\_hash} & \code{str}      & SHA-256 of the resulting VarFrame state.\\
\code{step}        & \code{PipelineStep} or \code{None} & Replayable call recorded for pipeline extraction (REQ-54): the VarFrame method, its keyword arguments, and the comment. \code{None} when the operation records none. Replay metadata, so it is excluded from the state hash.\\
\bottomrule
\end{tabularx}
\end{table}

History entries are indexed so that subranges can be exported as a Pipeline
(e.g.\ \code{db.history.to\_pipeline(start=2, end=7)}). The operations
between two VarFrame objects can be diffed by comparing their histories.

\subsection{State Hash Scope}
\label{sec:state-hash}

The reproducibility principle (P-03) is only testable if the state hash has
a defined scope. The \code{state\_hash} of a VarFrame covers: dimension
names and coordinate arrays; variable names and value arrays; the ordered
operation sequence with normalized arguments and comments; and the
uncertainty, correlation, and origin-tag arrays when present. It excludes
every volatile field: timestamps, user identity, source file paths, and the
\itc{} version, all of which live in Provenance (source content is covered
separately by \code{source\_hash}). Two runs that load the same bytes and
apply the same operations therefore produce identical state hashes on the
same platform and NumPy version. Across platforms, bitwise floating-point
equality is not guaranteed; cross-platform comparisons use
\code{pproc.compare} with tolerances, never hash equality (REQ-103).

\begin{ddbox}{DD-01}{Why split Provenance from History}
The split aligns \itc{} with the W3C PROV-DM model~\citep{provdm}, where
entity (the data), activity (the operation), and agent (the user) are
distinguishable concepts. It also enables serialization of provenance
graphs across multiple VarFrame objects, which arise naturally in
\code{pproc.compare(case\_a=db1, case\_b=db2, reference=db\_ref)} and in
\code{itc.aerospace} analyses that consume multiple inputs.
\end{ddbox}

\section{Axes and Custom Coordinate Frames}
\label{sec:axes}

\itc{} treats coordinate frames as first-class objects rather than implicit
context. The canonical reference frame is the AIAA body axis system per
AIAA R-004A-1992~\citep{aiaaR004A}. Additional named frames (stability,
wind, tunnel, propeller-disk, custom rig) are registered as \code{Axis}
objects.

\begin{lstlisting}
@dataclass(frozen=True)
class Axis:
    name: str                              # "body", "stability", "wind", ...
    rotation_matrix: np.ndarray | None     # 3x3, from canonical body axis
    angles_from: tuple[str, ...] | None = None  # condition-dependent frames
    parent: "Axis | None" = None           # for chained transforms
    description: str | None = None
\end{lstlisting}

A VarFrame carries an \code{AxisRegistry} that maps frame names to
\code{Axis} objects. The operation \code{db.rotate(target\_axis)}:

\begin{itemize}[noitemsep,topsep=2pt]
  \item detects vector groups by naming convention: triplets such as
        \code{(FX, FY, FZ)}, \code{(MX, MY, MZ)}, or any group declared
        explicitly via \code{db.declare\_vector("name", components=[...])};
  \item applies the rotation matrix to each detected group;
  \item propagates uncertainty using the rotation matrix as the Jacobian,
        which is exact since the rotation is linear, including covariance
        terms when a correlation matrix is present;
  \item updates the History with the source and target frames;
  \item preserves all HistoryFrame tags unchanged (rotation does not change
        whether a point is original, interpolated, or extrapolated).
\end{itemize}

This design supports the wind tunnel bookkeeping workflow, where balance
data routinely transit between rig, tunnel, body, stability, and wind axis
systems and where each transformation must remain auditable.

\subsection{Condition-Dependent Frames}
\label{sec:condition-frames}

A constant rotation matrix cannot express the frames that matter most in
practice: the transformation from body to wind or stability axes depends on
$\alpha$ and $\beta$, which are dimensions or variables of the VarFrame
itself, so the rotation varies from grid point to grid point. An
\code{Axis} therefore admits two definitions: a constant
\code{rotation\_matrix}, or a parametric definition via
\code{angles\_from=("alpha", "beta")} with a named angle convention, in
which case the matrix is evaluated per grid point. The rotation remains
linear in the vector components at every point, so uncertainty propagation
stays exact; when the referenced angles themselves carry uncertainty, the
sensitivity terms $\partial R / \partial \alpha$ enter the propagation by
the chain rule (REQ-101). The built-in \code{"wind"} and
\code{"stability"} frames are predefined this way per AIAA
R-004A-1992~\citep{aiaaR004A}.

\subsection{Moment Reference-Point Transfer}
\label{sec:moment-transfer}

Rotation alone does not close the balance bookkeeping chain: moments must
also transfer between reference points, from the balance moment center to
the model reference point or the aircraft center of gravity. The operation
\code{db.translate\_moments(to\_point=...)} applies
$\mathbf{M}' = \mathbf{M} + \mathbf{r} \times \mathbf{F}$ to each declared
force and moment vector group, where $\mathbf{r}$ is the offset between
reference points expressed in the current frame. The transfer is linear in
the force components, so the Jacobian is exact and covariance between force
and moment channels propagates correctly. The History entry records the
source point, the target point, and the frame in which the offset was
expressed (REQ-100).

\section{The .itceq Equation File}
\label{sec:itceq}

The \code{.itceq} file is the primary mechanism for reusable,
version-controlled analysis workflows. It is a TOML-structured text file
with five sections:

\begin{lstlisting}[language=toml]
[meta]
name        = "Aircraft balance campaign: power off"
version     = "1.0"
author      = "Geovana Neves"
description = "6-component internal balance, AIAA S-071A-1999 aligned"

[constants]
S_ref = 0.1963   # m^2, reference wing area
c_ref = 0.2526   # m, reference chord
b_ref = 0.800    # m, reference span

[uncertainties]
FX = 0.005       # N, balance calibration
FZ = 0.005       # N
MY = 0.0002      # N.m
V  = 0.02        # m/s, Pitot probe
rho = "0.05%"    # relative

[equations]
q_inf = "0.5 * rho * V**2"
CL    = "FZ / (q_inf * S_ref)"
CD    = "FX / (q_inf * S_ref)"
CM    = "MY / (q_inf * S_ref * c_ref)"

[corrections]
blockage = "1 + 0.005 * CL**2"
CL_corr  = "CL * blockage"
\end{lstlisting}

\subsection{Section Rules}

\begin{itemize}[noitemsep,topsep=2pt]
  \item \code{[meta]}: required. All sub-fields are optional strings.
  \item \code{[constants]}: key-value pairs, all numeric. Constants are
    registered before any equation is evaluated.
  \item \code{[uncertainties]}: values are float (absolute) or string
    ending in \code{"\%"} (relative). Per-variable.
  \item \code{[equations]}: evaluated in dependency order (topologically
    sorted at parse time when the optional \code{auto\_sort=True} flag is
    set, with a feedback message reporting the resolved order; otherwise
    evaluated in file order). Cyclic dependencies raise
    \code{ItceqCycleError} at parse time.
  \item \code{[corrections]}: evaluated after \code{[equations]}. May
    replace existing variables.
\end{itemize}

\section{The .itc Native Binary Format}

The \code{.itc} format is a ZIP archive containing the files listed in
Table~\ref{tab:itc-archive}. The format is designed for long-term archival:
it uses open standards (NumPy \code{.npz} and JSON) and can be inspected
without \itc{} using any ZIP tool plus \code{numpy.load}.

\begin{table}[H]
\centering
\caption{Layout of a \code{.itc} archive.}
\label{tab:itc-archive}
\small
\begin{tabularx}{\textwidth}{@{}lX@{}}
\toprule
\textbf{File inside archive} & \textbf{Content}\\
\midrule
\code{varframe.npz}      & All variable arrays in NumPy compressed format.\\
\code{dims.json}         & Dimension names, coordinate arrays, units, descriptions, numeric flags.\\
\code{vars\_meta.json}   & Variable units, descriptions, long names.\\
\code{uncframe.npz}      & Uncertainty arrays (absent when no UncFrame).\\
\code{correlation.json}  & Declared correlation matrix (absent when independent).\\
\code{historyframe.npz}  & Origin tags per variable (absent when all values are original).\\
\code{axes.json}         & Registered axes with rotation matrices (absent when only the canonical body axis is present).\\
\code{provenance.json}   & Origin record (Section~\ref{sec:provenance-history}).\\
\code{history.json}      & Ordered list of operation entries with indices.\\
\code{metadata.json}     & \itc{} version, format version, schema validation hash.\\
\bottomrule
\end{tabularx}
\end{table}

\section{The .itc\_pipe Pipeline Format}

A \code{.itc\_pipe} file is a human-readable serialization (JSON) of a
sequence of operations extracted from a VarFrame's History. DD-28
records why the encoding is JSON rather than the TOML this section named
at the 0.1.0 baseline: no Python version ships a standard-library TOML
writer, TOML has no null type (and \code{compute(fill=None)} is
meaningful, differing from the default \code{fill=nan}), and replay
arguments nest. The file records:

\begin{itemize}[noitemsep,topsep=2pt]
  \item the \itc{} version that created it,
  \item the index range in the source history,
  \item each operation's call with its keyword arguments and any
    attached comment (REQ-19),
  \item a content hash for integrity verification, reverified on load so
    a recipe edited after it was written fails loud rather than
    replaying something unintended.
\end{itemize}

Each recorded call is validated against the set of replayable operations
when the file is read. A \code{.itc\_pipe} is meant to be shared and
version-controlled, so a hand-edited file naming an arbitrary method is
refused rather than dispatched.

A \code{.itc\_pipe} file is intentionally separate from \code{.itc}: a
pipeline encapsulates how to transform data without containing data itself,
so it can be shared, reviewed, and version-controlled alongside the source
code.

\section{The .itc\_surr Surrogate Format}

A \code{.itc\_surr} file is a binary serialization of a fitted surrogate
model with associated metadata. It is exportable independently of any
VarFrame, enabling reuse of trained surrogates in performance and EUB
calculations without reloading the source dataset.

By default the file stores only the model state plus the SHA-256 hash of
the source VarFrame; passing \code{embed\_source=True} to
\code{surr.save(...)} additionally embeds the full source VarFrame in the
archive, yielding a self-contained but larger file. The trade-off is
chosen explicitly per save.
